\documentclass[border=4pt]{standalone}
\usepackage{amsmath,amssymb,mathtools,xcolor,varwidth}
\begin{document}
\begin{varwidth}{28em}
\large\bfseries
Pick any one pair of homologous sides: $AB$ in the first figure answers to $DE$ in the second. Newton's claim is that \emph{every} homologous pair stands in one and the same ratio --- so $AB/DE = BC/EF = CA/FD = k$, a single number $k$ that measures how much larger the one figure is than the other. But the \emph{areas} do not grow in this same ratio: they grow in its \emph{duplicate ratio}, that is, in the ratio $k^2$. Thus if the larger figure's sides are twice those of the smaller ($k = 2$), its area is not twice but \emph{four} times as great ($k^2 = 4$); triple the side, and the area is ninefold. This is why doubling a map's scale quadruples the ink, and why the homologous sides --- highlighted here --- govern both lengths and areas at once. \textit{Q.E.D.}
\end{varwidth}
\end{document}
